\section{Background and Related Work}
\label{sec:background}

Our work sits at the intersection of Retrieval-Augmented Generation (RAG), autonomous security agents, and structural information retrieval. This section reviews the foundational literature and recent advancements that inform the design of VulnRAG.

\subsection{Retrieval-Augmented Generation (RAG)}
The paradigm of augmenting Large Language Models with external knowledge was formalized by Lewis et al. \cite{rag}, who combined a dense retriever with a sequence-to-sequence generator. Guu et al. \cite{realm} further emphasized that external knowledge corpora (e.g., Wikipedia) are more easily updated than parametric model weights, a critical requirement for the rapidly evolving CVE landscape. However, Liu et al. \cite{lostmiddle} identified the ``Lost in the Middle'' problem, where LLMs fail to utilize information placed in the center of long contexts. While general reranking has been suggested as a mitigation, VulnRAG provides a domain-specific implementation through Multi-Factor Semantic Reranking (MFSR) that prioritizes operational actionability over pure text similarity.

\subsection{Agentic Workflows and Self-Reflection}
The transition from static to agentic RAG is driven by the need for iterative reasoning and self-correction. Yao et al. introduced \textbf{ReAct} \cite{react} and \textbf{Tree of Thoughts (ToT)} \cite{tot}, enabling models to explore multiple strategic branches. Shinn et al.'s \textbf{Reflexion} \cite{reflexion} demonstrated that verbal reinforcement learning can improve reasoning performance without fine-tuning. Most relevantly, \textbf{Self-RAG} \cite{selfrag} and \textbf{ActiveRAG} \cite{activerag} proposed models that adaptively retrieve and critique their own findings. However, Self-RAG requires specialized reflection tokens and model training. In contrast, VulnRAG implements an agentic self-correction loop zero-shot, utilizing a decoupled architecture of specialized agents to optimize retrieval quality in a provider-agnostic manner.

\subsection{Autonomous Security Intelligence}
The application of LLMs to cybersecurity has led to the development of sophisticated pentesting agents. \textbf{PentestGPT} \cite{pentestgpt} introduced the Pentesting Task Tree (PTT) to maintain state during multi-step assessments. Recent works like \textbf{RAVEN} \cite{raven} and \textbf{AutoAttacker} \cite{autoattacker} have shown success in generating vulnerability reports and reproducing exploits. Despite these gains, existing systems often treat vulnerability intelligence as a ``black box'' tool output. Multi-level RAG models \cite{multilevelrag} have begun integrating multiple security databases (NVD, CERT, GitHub), but they primarily focus on the data ingestion layer. VulnRAG fills this gap by focusing on the \textit{agentic intelligence layer}—implementing structural discovery and complexity-aware routing that SOTA security systems currently lack.

\subsection{Structural and Graph-based Retrieval}
Traditional RAG relies on flat chunking and vector similarity. To overcome this, \textbf{RAPTOR} \cite{raptor} proposed recursive abstractive processing to build a tree of document summaries. While effective for thematic retrieval, RAPTOR does not capture the explicit structural relationships found in vulnerability databases. VulnRAG leverages **CVE Relationship Graphs**—where nodes are vulnerabilities and edges represent explicit reference URLs or product clusters—to discover attack chains. This moves the retrieval frontier from "Finding similar text" to "Discovering linked threats," providing a unique discovery yield for security analysts.
